\rtmtight
\begin{tblr}{width=\textwidth,colspec={|Q[c,m,2.1cm]|X[l,m]|},hlines,vlines,hspan=minimal,rowsep=1.5pt,colsep=3pt,row{1}={bg=headgray,font=\bfseries}}
\SetCell{c,m}\hfil\logo\hfil & \textbf{POLITEKNIK NEGERI MALANG}\par\textbf{JURUSAN TEKNOLOGI INFORMASI}\par\textbf{PROGRAM STUDI: D4 TEKNIK INFORMATIKA} \\
\SetCell[c=2]{l} \textbf{RENCANA TUGAS MAHASISWA (RTM) DAN RUBRIK PENILAIAN} & \\
Nama Mata Kuliah & Komunikasi dan Etika Profesi \\
Kode Mata Kuliah & RTI256001 \quad \textbf{sks / Jam perminggu:} 2 SKS/4 jam \quad \textbf{Semester:} 6 \\
Dosen Pengampu & Tim Pengajar Program Studi D4 TI \\
\end{tblr}
\assessmentblock{Tugas Etika Profesi dan Komunikasi}{Case Method dan studi kasus}{SCPMK0101-04501, SCPMK0301-04502, SCPMK0303-04503}{Mahasiswa menyelesaikan tugas menganalisis kasus etika profesi dan hukum ITE, mempraktikkan komunikasi profesional, dan merumuskan rencana awal pengembangan karir.}{Menganalisis kasus, menyusun rekomendasi etis, membuat artefak komunikasi profesional, mendokumentasikan, dan mempresentasikan.}{Laporan analisis kasus, dokumen komunikasi profesional, rencana karir awal, dan/atau bahan presentasi.}{Analisis kasus etika profesi dan hukum ITE & 14 & Kasus pelanggaran etika, hukum ITE, dan K3 dianalisis secara mendalam, rekomendasi tindakan tepat dan berbasis regulasi yang berlaku, serta argumen etis dirumuskan dengan jelas. & Analisis kasus cukup baik tetapi rekomendasi atau argumen etis belum sepenuhnya berbasis regulasi yang tepat. & Analisis tidak mendalam, rekomendasi tidak realistis, atau tidak mempertimbangkan regulasi yang relevan. \\ \\
Kualitas komunikasi dan presentasi teknis & 10 & Komunikasi profesional (email, laporan, presentasi) memenuhi standar korporat, pesan terstruktur, bahasa sesuai konteks, dan disampaikan dengan percaya diri. & Komunikasi profesional cukup baik tetapi struktur, bahasa, atau penyampaian masih perlu ditingkatkan. & Komunikasi tidak memenuhi standar profesional atau tidak dapat menyampaikan pesan secara efektif. \\ \\
Perencanaan awal pengembangan karir & 6 & Rencana karir awal dirumuskan dengan tujuan yang spesifik, kompetensi yang diperlukan diidentifikasi, dan langkah konkret disusun berbasis kondisi industri TI. & Rencana karir ada tetapi tujuan, kompetensi, atau langkah konkret belum sepenuhnya spesifik. & Rencana karir tidak realistis atau tidak berbasis kompetensi TI yang relevan. \\}

\assessmentblock{Kuis 1 Etika dan Hukum TI}{Kuis}{SCPMK0101-04501}{Kuis tertulis untuk mengukur pemahaman etika profesi TI, kode etik ACM/IEEE, hukum ITE, dan K3.}{Menjawab soal pilihan ganda dan/atau uraian terkait etika profesi, regulasi hukum, dan K3.}{Jawaban tertulis kuis.}{Pemahaman etika profesi dan kode etik ACM/IEEE & 3 & Prinsip kode etik ACM/IEEE, nilai profesionalisme, dan aplikasinya pada skenario TI dijelaskan dengan benar dan contoh penerapannya tepat. & Sebagian besar konsep dipahami tetapi penerapan pada skenario spesifik masih kurang lengkap. & Kode etik tidak dipahami atau contoh penerapan keliru. \\ \\
Analisis regulasi hukum ITE dan K3 & 2 & Pasal-pasal hukum ITE yang relevan dan regulasi K3 di lingkungan TI diidentifikasi dengan benar dan implikasinya pada kasus dianalisis tepat. & Analisis regulasi cukup baik tetapi beberapa pasal atau implikasi masih belum tepat. & Analisis tidak tepat atau regulasi yang digunakan tidak relevan. \\}

\assessmentblock{UTS Kasus Etika, Hukum, dan Komunikasi}{UTS berbasis kasus}{SCPMK0101-04501, SCPMK0301-04502}{Evaluasi tengah semester untuk mengukur kemampuan menganalisis kasus etika profesi, hukum ITE, dan berkomunikasi secara profesional.}{Menganalisis kasus komprehensif terkait etika profesi dan membuat artefak komunikasi profesional.}{Dokumen analisis kasus dan artefak komunikasi profesional.}{Analisis dan argumentasi kasus etika profesi & 7 & Kasus dianalisis secara mendalam menggunakan kerangka etika dan regulasi yang tepat, rekomendasi tindakan dapat dipertahankan dengan argumen yang kuat dan berbasis bukti. & Analisis cukup baik tetapi argumentasi atau penggunaan kerangka etika masih belum kuat. & Analisis tidak mendalam atau tidak dapat memberikan rekomendasi yang dapat dipertahankan. \\ \\
Kualitas komunikasi profesional tertulis dan lisan & 6 & Artefak komunikasi (laporan/presentasi) memenuhi standar profesional, pesan terstruktur, bahasa tepat, dan disampaikan secara efektif. & Komunikasi cukup profesional tetapi struktur, bahasa, atau penyampaian masih perlu ditingkatkan. & Komunikasi tidak memenuhi standar profesional atau tidak dapat menyampaikan pesan secara efektif. \\}

\assessmentblock{Kuis 2 Isu Sosial dan Karir TI}{Kuis}{SCPMK0303-04503}{Kuis tertulis untuk mengukur pemahaman isu sosial teknologi dan strategi perencanaan karir di bidang TI.}{Menjawab soal terkait isu privasi data, tanggung jawab sosial korporasi, personal branding, dan perencanaan karir.}{Jawaban tertulis kuis.}{Pemahaman isu sosial teknologi dan tanggung jawab korporasi & 3 & Isu privasi data, surveillance, dan tanggung jawab sosial korporasi TI dijelaskan dengan benar dan contoh kasus aktual dianalisis tepat. & Sebagian besar isu dipahami tetapi analisis kasus atau implikasi masih kurang lengkap. & Isu sosial tidak dipahami atau analisis tidak relevan. \\ \\
Strategi personal branding dan perencanaan karir & 2 & Komponen personal branding (LinkedIn, GitHub, portofolio) dan strategi perencanaan karir di industri TI dijelaskan dengan tepat dan aplikatif. & Strategi sebagian besar tepat tetapi beberapa komponen atau aplikasinya belum lengkap. & Strategi tidak tepat atau tidak relevan dengan industri TI. \\}

\assessmentblock{PBL Simulasi Wawancara dan Presentasi Portofolio}{Project Based Learning}{SCPMK0301-04502, SCPMK0303-04503}{Mahasiswa melakukan simulasi wawancara kerja profesional, mempresentasikan portofolio digital, dan menyajikan rencana pengembangan karir berbasis kompetensi TI.}{Mempersiapkan portofolio, melakukan simulasi wawancara, mempresentasikan rencana karir, dan menerima umpan balik.}{Portofolio digital, bahan presentasi, rekaman simulasi wawancara, dan rencana karir tertulis.}{Kualitas simulasi wawancara dan presentasi teknis & 10 & Simulasi wawancara menunjukkan komunikasi profesional yang efektif, jawaban terstruktur dengan metode STAR, dan presentasi teknis disampaikan dengan percaya diri serta meyakinkan. & Simulasi wawancara dan presentasi cukup baik tetapi beberapa aspek komunikasi atau struktur jawaban masih perlu ditingkatkan. & Simulasi wawancara tidak menunjukkan komunikasi profesional atau presentasi tidak dapat menyampaikan pesan secara efektif. \\ \\
Kelengkapan portofolio profesional & 7 & Portofolio digital lengkap mencakup profil profesional, proyek, sertifikasi, dan artefak kerja yang relevan, serta dipresentasikan secara kohesif. & Portofolio sebagian besar lengkap tetapi beberapa komponen atau kohesivitas presentasi masih belum optimal. & Portofolio tidak lengkap atau tidak mencerminkan kompetensi TI yang relevan. \\ \\
Rencana pengembangan karir berbasis kompetensi & 5 & Rencana karir tertulis dengan tujuan SMART, peta kompetensi yang diperlukan, dan langkah konkret berbasis kondisi industri TI yang realistis. & Rencana karir ada tetapi tujuan, peta kompetensi, atau langkah konkret belum sepenuhnya SMART dan realistis. & Rencana karir tidak terstruktur atau tidak realistis untuk kondisi industri TI. \\}

\assessmentblock{UAS Presentasi Portofolio dan Rencana Karir}{UAS presentasi dan defense}{SCPMK0303-04503}{Mahasiswa mempresentasikan portofolio profesional final dan rencana pengembangan karir secara komprehensif, mempertahankan keputusan berdasarkan kompetensi dan tren industri TI.}{Mempresentasikan portofolio, menjelaskan rencana karir, dan menjawab pertanyaan dari penguji.}{Portofolio digital final, rencana karir komprehensif, dan bahan presentasi.}{Kelengkapan dan kualitas portofolio digital & 10 & Portofolio digital mencakup semua komponen yang diperlukan (profil, proyek, sertifikasi, artefak kerja), terpresentasi secara profesional, dan mencerminkan kompetensi TI yang relevan. & Portofolio sebagian besar lengkap dan profesional tetapi beberapa komponen atau kualitas presentasinya masih perlu ditingkatkan. & Portofolio tidak lengkap atau tidak mencerminkan standar profesional. \\ \\
Rencana pengembangan karir berbasis kompetensi & 9 & Rencana karir komprehensif dengan analisis kesenjangan kompetensi, tujuan jangka pendek-menengah-panjang yang SMART, dan strategi pencapaian yang realistis berbasis tren industri TI. & Rencana karir cukup baik tetapi analisis kesenjangan, tujuan, atau strategi pencapaian masih belum sepenuhnya komprehensif. & Rencana karir tidak komprehensif atau tidak berbasis analisis kompetensi yang realistis. \\ \\
Presentasi dan defense portofolio & 6 & Presentasi disampaikan dengan percaya diri dan terstruktur, keputusan pengembangan karir dapat dipertahankan dengan argumen yang kuat, dan pertanyaan dijawab secara meyakinkan. & Presentasi cukup baik tetapi beberapa argumen atau jawaban atas pertanyaan masih belum kuat. & Tidak dapat mempresentasikan portofolio secara profesional atau menjawab pertanyaan dengan tepat. \\}

\begin{tblr}{width=\textwidth,colspec={|Q[l,m,3.2cm]|X[l,m]|},hlines,vlines,hspan=minimal,row{1-3}={bg=headgray},rowsep=1.5pt,colsep=3pt}
Jadwal Pelaksanaan: & Minggu 1--17 sesuai RPS. Setiap evaluasi memiliki RTM dan rubrik multi-kriteria. \\
Lain-Lain Yang Diperlukan: & LMS, dokumen portofolio, perangkat lunak sesuai mata kuliah, dan perangkat presentasi. \\
Pustaka: & Mengacu pustaka utama dan pendukung pada bagian RPS. \\
\end{tblr}
